% !TeX root = ../main.tex
% ============================================================
%  Capítulo 3 – Materiais e Métodos
% ============================================================
\chapter{Materiais e Métodos}
\label{cap:materiais_metodos}

% TODO: Descrever a metodologia completa: dados, modelos, avaliação.
%       Deve ser suficientemente detalhado para permitir reprodutibilidade.

% ============================================================
\section{Dados Meteorológicos e Variáveis de Entrada}
\label{sec:dados}

% ------------------------------------------------------------
\subsection{O conjunto de dados das usinas}
\label{subsec:dataset_usinas}

% TODO: Descrever as usinas fotovoltaicas utilizadas: localização, capacidade
%       instalada, período de dados, resolução temporal.
%       Indicar a fonte (operador, SCADA, ONS, etc.) e eventuais acordos de
%       confidencialidade.

% ------------------------------------------------------------
\subsection{Variáveis calculadas}
\label{subsec:variaveis_calculadas}

% TODO: Detalhar variáveis derivadas: ângulo solar (elevação, azimute),
%       índice clear-sky (GHI medido / GHI clear-sky), temperatura de célula
%       estimada etc. Citar as equações utilizadas.

% ------------------------------------------------------------
\subsection{Pré-processamento dos dados}
\label{subsec:preprocessamento}

% TODO: Descrever tratamento de outliers, dados faltantes (interpolação,
%       exclusão), remoção de dados noturnos e filtros de qualidade.
%       Quantificar a taxa de dados ausentes ou descartados.

% ------------------------------------------------------------
\subsection{Engenharia de \textit{features}}
\label{subsec:feature_engineering}

% TODO: Explicar a criação de lags temporais, médias móveis, codificação
%       cíclica de variáveis de tempo (hora, dia do ano), interações entre
%       variáveis. Justificar cada decisão com base na física do problema.

% ------------------------------------------------------------
\subsection{Conjunto final de variáveis}
\label{subsec:conjunto_variaveis}

% TODO: Apresentar tabela-resumo com todas as variáveis de entrada finais:
%       nome, unidade, descrição, origem (medida / calculada / NWP).

% Exemplo de tabela:
% \begin{table}[htb]
%   \centering
%   \caption{Variáveis de entrada utilizadas nos modelos.}
%   \label{tab:variaveis_entrada}
%   \begin{tabularx}{\textwidth}{llXl}
%     \toprule
%     \textbf{Variável} & \textbf{Unidade} & \textbf{Descrição} & \textbf{Origem} \\
%     \midrule
%     GHI               & \si{\watt\per\square\meter} & Irradiância horizontal global & Medida \\
%     $T_{amb}$         & \si{\celsius} & Temperatura do ar & Medida \\
%     ...               & ...           & ...               & ...    \\
%     \bottomrule
%   \end{tabularx}
%   \fonte{Elaborado pelo autor.}
% \end{table}

% ============================================================
\section{Modelos Adotados}
\label{sec:modelos_adotados}

% TODO: Descrever cada modelo selecionado, p.ex.:
%   \subsection{Gradient Boosting (XGBoost / LightGBM)}
%   \subsection{Random Forest}
%   \subsection{Redes Neurais LSTM}
%   Para cada um: princípio de funcionamento, hiperparâmetros principais,
%   implementação utilizada (biblioteca Python/scikit-learn/etc.).

% ============================================================
\section{Métricas de Avaliação e \textit{Benchmarks}}
\label{sec:metricas}

% TODO: Definir e justificar as métricas escolhidas. Exemplos:
%
% \textbf{RMSE} (Root Mean Square Error):
% \begin{equation}
%   \text{RMSE} = \sqrt{\frac{1}{n}\sum_{i=1}^{n}(\hat{y}_i - y_i)^2}
%   \label{eq:rmse}
% \end{equation}
%
% \textbf{MAE} (Mean Absolute Error):
% \begin{equation}
%   \text{MAE} = \frac{1}{n}\sum_{i=1}^{n}|\hat{y}_i - y_i|
%   \label{eq:mae}
% \end{equation}
%
% \textbf{MAPE} (Mean Absolute Percentage Error):
% \begin{equation}
%   \text{MAPE} = \frac{100\%}{n}\sum_{i=1}^{n}\left|\frac{\hat{y}_i - y_i}{y_i}\right|
%   \label{eq:mape}
% \end{equation}
%
% Discutir limitações de cada métrica (e.g., sensibilidade do MAPE a y ≈ 0).
% Descrever como os benchmarks (baselines) serão comparados aos modelos propostos.
